\documentclass[12pt,a4paper]{article}

% ── 中文支持 ──
\usepackage[UTF8]{ctex}
\usepackage{fontspec}
\setmainfont{Times New Roman}
\setCJKmainfont{SimSun}[AutoFakeBold]

% ── 页面布局 ──
\usepackage[margin=2.5cm]{geometry}
\usepackage{graphicx}
\usepackage{float}
\usepackage{amsmath,amssymb}
\usepackage{booktabs}
\usepackage{hyperref}
\usepackage{enumitem}
\usepackage{listings}
\usepackage{xcolor}

\hypersetup{
    colorlinks=true,
    linkcolor=blue,
    citecolor=blue,
    urlcolor=blue,
}

\lstset{
    basicstyle=\ttfamily\small,
    breaklines=true,
    frame=single,
    numbers=left,
    numberstyle=\tiny,
    numbersep=5pt,
    backgroundcolor=\color{gray!5},
}

\title{\bfseries 小样本条件下中文多角色文本风格迁移研究}
\author{
    王浩楠 \quad 张博凇 \quad 范智杰 \quad 何顺康 \quad 陈佳铭 \\[4pt]
    \small 机器学习课程项目组
}
\date{\today}

\begin{document}

\maketitle

\begin{abstract}
本文研究小样本条件下的中文多角色文本风格迁移任务。我们构建了一个包含 5 类角色关系（老板、同事、好朋友、女朋友、母亲）的中文平行语料库，每条中性文本对应 5 种不同角色风格的改写。基于该数据集，我们提出了一种多 LoRA 适配器架构：以 mT5-large 为基座模型，为每个目标角色独立训练一个 LoRA 适配器，实现角色感知的语用生成。此外，我们开发了一个完整的 Web 演示系统，支持教师模型（DeepSeek API）与微调模型的在线对比。实验表明，在约 8000 条训练样本下，多 LoRA 方案能够有效捕捉不同关系角色的语用差异。
\end{abstract}

\section{引言}

在日常交流中，同一语义内容面对不同关系角色时，语气、称呼、礼貌程度和情绪表达存在显著差异。例如，告知“今天加班”这一信息，对老板应当正式汇报，对女朋友则应表达在意和安抚。现有的文本风格迁移研究较少关注关系角色维度，本文旨在填补这一空白。

\section{任务定义}

\begin{itemize}[leftmargin=2em]
    \item \textbf{输入}：中性文本 + 目标角色（5 选 1）
    \item \textbf{输出}：符合该角色关系的改写文本
    \item \textbf{核心挑战}：模型需要理解“对谁说话”，并据此调整称呼、礼貌层级、信息组织方式
\end{itemize}

5 个目标角色及其语用特征见表~\ref{tab:roles}。

\begin{table}[H]
\centering
\caption{角色定义与语用特征}
\label{tab:roles}
\begin{tabular}{c|c|c|c}
\toprule
角色 & 权力 & 距离 & 核心要求 \\
\midrule
老板   & 上→下 & 远 & 正式、汇报、对结果负责 \\
同事   & 平等   & 中 & 协作、留有余地、不卑不亢 \\
好朋友 & 平等   & 近 & 随意自然、直来直去、可吐槽 \\
女朋友 & 平等   & 最近 & 亲密温柔、情绪优先、安抚 \\
母亲   & 下→上 & 近 & 尊重温暖、让长辈放心 \\
\bottomrule
\end{tabular}
\end{table}

\section{数据构建}

\subsection{数据来源与筛选}

我们从 LCCC\cite{wang2020chinese}（Large-scale Cleaned Chinese Conversation）数据集中抽取 15,000 条对话，经过以下管线处理：

\begin{enumerate}[leftmargin=2em]
    \item \textbf{拆句}：将多轮对话按标点拆分为单句
    \item \textbf{多阶段过滤}：长度（8--65 字）、第一人称、排除角色标记与性别暴露、信息量筛选
    \item \textbf{LLM 批量改写}：调用 DeepSeek API，将中性文本并行生成 5 类角色改写
    \item \textbf{质量校验}：JSON 结构校验、角色区分度检查（ROUGE-L）、语义适配性判断
    \item \textbf{格式转换}：展开为“中性句 $\rightarrow$ 角色改写”的平行语料
\end{enumerate}

\subsection{数据统计}

最终数据集包含约 2,000 条中性文本，展开为约 8,000 条训练样本（过滤掉 N/A 不适配的条目），按 8:1:1 划分训练/验证/测试集。各角色样本分布见表~\ref{tab:data}。

\begin{table}[H]
\centering
\caption{各角色训练样本数}
\label{tab:data}
\begin{tabular}{c|c}
\toprule
角色 & 样本数 \\
\midrule
老板       & $\sim$1,500 \\
同事       & $\sim$1,500 \\
好朋友     & $\sim$1,600 \\
女朋友     & $\sim$1,600 \\
母亲       & $\sim$1,500 \\
\bottomrule
\end{tabular}
\end{table}

\section{模型微调}

\subsection{基座模型选择}

我们选用 mT5-large\cite{xue2021mt5}（1.2B 参数）作为基座模型。mT5 基于 T5 的 encoder-decoder 架构，在大规模多语言语料上预训练，具备较好的跨语言迁移能力。与中文原生 T5 相比，mT5-large 参数量更大、语言知识更丰富，更适合捕捉角色风格的细腻差异。

\subsection{多 LoRA 适配器架构}

传统的 LoRA 微调在整个任务数据上训练单一适配器，不同角色共享同一组低秩矩阵。这种方式的缺陷在于角色间的风格信号容易相互干扰。我们提出多 LoRA 适配器方案：为每个目标角色独立训练一个 LoRA 适配器，基座模型参数完全冻结。架构如图~\ref{fig:architecture} 所示。

\begin{figure}[H]
\centering
% 用文字描述替代图片
\begin{verbatim}
           ┌──────────────────────┐
           │   mT5-large (冻结)    │
           └──┬──┬──┬──┬──┬───────┘
              │  │  │  │  │
          LoRA LoRA LoRA LoRA LoRA
          老板  同事 好朋友 女友  母亲
\end{verbatim}
\caption{多 LoRA 适配器架构}
\label{fig:architecture}
\end{figure}

推理时，根据用户选择的目标角色，动态切换对应的 LoRA 适配器：

\begin{lstlisting}[language=Python]
# 加载基座 + 全部适配器
base_model = AutoModelForSeq2SeqLM.from_pretrained("mT5-large")
model = PeftModel.from_pretrained(base_model, "output/boss/final",
                                   adapter_name="boss")
model.load_adapter("output/colleague/final", adapter_name="colleague")
# ... 加载其余角色

# 推理时切换
model.set_adapter("boss")
output = model.generate(**inputs)
\end{lstlisting}

\subsection{训练配置}

\begin{table}[H]
\centering
\caption{LoRA 训练超参数}
\label{tab:hyperparams}
\begin{tabular}{c|c}
\toprule
超参数 & 值 \\
\midrule
LoRA rank $r$      & 16 \\
LoRA alpha         & 32 \\
LoRA dropout       & 0.1 \\
目标模块            & q, v (注意力层) \\
优化器              & AdamW \\
学习率              & $1\times10^{-4}$ \\
学习率调度          & warmup 10\%, 线性衰减 \\
Batch size         & 4 \\
最大 Epoch          & 8 \\
早停耐心值          & 5 \\
梯度裁剪            & 1.0 \\
精度                & float32 \\
可训参数            & 2,359,296 (0.19\%) \\
\bottomrule
\end{tabular}
\end{table}

每个角色的 LoRA 适配器仅包含约 236 万可训练参数（基座模型的 0.19\%），训练耗时约 30 分钟/角色（单卡 RTX 4080），5 个角色总计约 2.5 小时。

\subsection{训练过程}

训练过程中观察到以下现象：

\begin{enumerate}[leftmargin=2em]
    \item \textbf{FP16 稳定性问题}：mT5-large 在 float16 精度下梯度频繁出现 NaN，切换为 float32 并添加梯度裁剪（max\_grad\_norm=1.0）后稳定。
    \item \textbf{早停触发}：eval loss 在约 3 轮后趋于平稳，EarlyStopping 自动终止训练，避免过拟合。
    \item \textbf{各角色收敛趋势一致}：5 个角色的 loss 曲线形态相似，最终 eval loss 均在 1.3--1.5 范围内。
\end{enumerate}

\section{系统实现}

\subsection{后端架构}

后端基于 FastAPI 构建，主要接口见表~\ref{tab:api}。

\begin{table}[H]
\centering
\caption{后端 API 接口}
\label{tab:api}
\begin{tabular}{c|c|c}
\toprule
端点 & 方法 & 功能 \\
\midrule
/api/auth/register & POST & 用户注册 \\
/api/auth/login    & POST & 登录，返回 JWT \\
/api/meta/roles    & GET  & 角色列表 \\
/api/meta/models   & GET  & 模型列表 \\
/api/transfers     & POST & 提交改写请求 \\
/api/transfers     & GET  & 历史记录 \\
\bottomrule
\end{tabular}
\end{table}

推理服务支持两种模型模式：

\begin{itemize}[leftmargin=2em]
    \item \textbf{deepseek\_api}：调用 DeepSeek Chat API，作为教师模型基线
    \item \textbf{lora\_finetuned}：加载本地 mT5-large + 多 LoRA 适配器，按角色切换
\end{itemize}

\subsection{前端设计}

前端基于 Vue 3 + Vite 构建，采用左右分栏布局：左侧栏展示历史记录与用户信息，右侧主区域为对话式交互界面。用户依次选择目标角色与模型类型，输入中性文本后获得改写结果。界面参考 DeepSeek 网页端设计风格，支持注册/登录、历史记录回溯。

\section{实验结果与分析}

\subsection{定量评估}

\begin{table}[H]
\centering
\caption{各角色测试集 eval loss}
\label{tab:results}
\begin{tabular}{c|c}
\toprule
角色 & Eval Loss \\
\midrule
老板       & 1.42 \\
同事       & 1.38 \\
好朋友     & 1.35 \\
女朋友     & 1.40 \\
母亲       & 1.39 \\
\bottomrule
\end{tabular}
\end{table}

各角色 eval loss 在 1.3--1.5 之间，差异较小，说明多 LoRA 方案在不同角色上均能有效学习。

\subsection{推理样例}

\begin{table}[H]
\centering
\caption{推理样例}
\label{tab:samples}
\small
\begin{tabular}{c|c|c}
\toprule
角色 & 输入 & 输出 \\
\midrule
母亲 & 我今天吃的都感觉很好吃 & 妈，我今天吃的都感觉很好吃 \\
同事 & 遇到喜欢的就大胆去追 & 遇到喜欢的就大胆去追，有空一起吃饭 \\
老板 & 吃东西好多了，只是缺少睡眠 & 吃东西好多了，只是缺少睡眠，白天忙店铺，晚上忙论文 \\
\bottomrule
\end{tabular}
\end{table}

从样例中可以观察到，模型学会了基本的称呼添加（如“妈”），但在语气精细度上仍有提升空间。部分输出接近原文直译，说明当前数据量和 LoRA rank 可能不足以完全捕捉角色间的语用差异。

\section{结论}

本文提出了一个多 LoRA 适配器架构用于中文多角色文本风格迁移，在约 8,000 条训练样本下实现了基本的角色风格改写能力。主要贡献包括：（1）构建了一个 5 角色中文平行语料库；（2）验证了多 LoRA 方案在角色风格迁移任务上的有效性；（3）开发了完整的 Web 演示系统。未来工作包括扩大数据集规模、引入更细粒度的语用标注、以及探索更大规模的中文预训练模型。

\bibliographystyle{plain}
\begin{thebibliography}{99}

\bibitem{wang2020chinese}
Y.~Wang, P.~Ke, Y.~Zheng, K.~Huang, Y.~Jiang, X.~Zhu, and M.~Huang,
\emph{A Large-Scale Chinese Short-Text Conversation Dataset},
NLPCC, 2020.

\bibitem{xue2021mt5}
L.~Xue, N.~Constant, A.~Roberts, M.~Kale, R.~Al-Rfou, A.~Siddhant, A.~Barua, and C.~Raffel,
\emph{mT5: A Massively Multilingual Pre-trained Text-to-Text Transformer},
NAACL, 2021.

\bibitem{hu2022lora}
E.~J.~Hu, Y.~Shen, P.~Wallis, Z.~Allen-Zhu, Y.~Li, S.~Wang, L.~Wang, and W.~Chen,
\emph{LoRA: Low-Rank Adaptation of Large Language Models},
ICLR, 2022.

\bibitem{deepseek2025}
DeepSeek,
\emph{DeepSeek Chat API Documentation},
2025.

\end{thebibliography}

\end{document}
